\section{Method}
\label{sec:method}

\subsection{Preliminaries and Notation}
\label{sec:prelim}

% Define only the notation the method section actually uses. A reviewer should
% be able to read this subsection once and never look back at it.
\ARTODO{problem setup and notation}

\subsection{\ARTODO{name of the method}}
\label{sec:method-core}

% State the method precisely enough that a competent reader could reimplement
% it from this subsection alone. Equations get numbers only if referenced.
\ARTODO{the core construction}

% \begin{equation}
%   \label{eq:main}
%   % the central equation, once the method is settled
% \end{equation}

\subsection{Why It Should Work}
\label{sec:intuition}

% The argument that connects the construction to the claimed effect. This can
% be a proof, a bound, or an honest intuition -- but it must be one of them,
% and if it is intuition, say so rather than dressing it as analysis.
\ARTODO{the mechanism, stated as a falsifiable claim}

\subsection{Cost}
\label{sec:cost}

% Compute, memory and wall-clock relative to the baseline. Reviewers ask for
% this whenever a method adds machinery; supplying it unprompted is cheap.
\ARTODO{asymptotic and measured overhead, using \ARnum{} until measured}
